% Copyright (c) 2026 Geovana Neves. Licensed under CC BY 4.0 (https://creativecommons.org/licenses/by/4.0/). See LICENSE-AND-AUTHORSHIP.md.
%
% 02-gui-to-pyfs/text/03-boundary-conditions.tex
%
% Transcribed from the tables of docs/gui-to-pyfs.md, one row per GUI step:
% the key that takes it (links and solver commands dropped) and the builds
% its commands are verified on. When the page changes a row, change the same
% row here, and the counts of 01-reading.tex with it.

\section{Boundary conditions}

\begin{frame}{Boundary conditions: what a key takes (1/3)}
\relax
{\ftstablesize\renewcommand{\arraystretch}{1.05}
\begin{tabularx}{\linewidth}{@{}JKZ@{}}
\guihead
Mark the trailing edges from a file of edge points & sidecar table \code{[trailing\_edges]} with \code{file}, and \code{type} and \code{tolerance} where not the default & 26.124 \\
Detect the trailing edges over the whole mesh & sidecar table \code{[trailing\_edges]} with \code{detect = "auto"} & 26.101 to 26.122 \\
Detect the trailing edges of the surfaces you name, above a sweep angle & sidecar table \code{[trailing\_edges]} with \code{detect = \{ surfaces = [...], sweep\_angle = ... \}} & none \\
Mark the wake-termination nodes, over the whole mesh or by surface & sidecar table \code{[wake\_termination]} & none \\
Make the boundaries you name base regions & row key \code{BASE\_REGIONS}, or pproc key \code{base\_regions}, naming the base and not the body that carries it & none \\
Detect the base regions over the whole mesh & sidecar table \code{[base\_regions]} with \code{detect = "auto"} & none \\
Apply or skip the trailing edges, wake-termination nodes and base regions the geometry declares & setup keys \code{apply\_trailing\_edges}, \code{apply\_wake\_termination} and \code{apply\_base\_regions}; \code{false} skips the redefinition and never clears saved state & none \\
\bottomrule
\end{tabularx}}
\guisource{Boundary conditions}
\end{frame}

\begin{frame}{Boundary conditions: what a key takes (2/3)}
\relax
{\ftstablesize\renewcommand{\arraystretch}{1.05}
\begin{tabularx}{\linewidth}{@{}JKZ@{}}
\guihead
Change one trailing edge's type: relaxed, jet outflow, vortex shedding & setup key \code{trailing\_edge\_types}, each trailing-edge index with its type & 26.121 to 26.124 \\
Mark wake-termination nodes by hand, or stop a trailing edge's wake & setup keys \code{mark\_wake\_termination\_nodes} and \code{disabled\_wake\_trailing\_edges} & 26.121 to 26.124, except one \\
Shed wakes from the leading edges of a surface & setup key \code{leading\_edge\_wake\_boundaries} & none \\
Create, remesh, delete or bend a base region, set its pressure, select its faces or mark its trailing or outflow edges & setup key \code{base\_region\_operations}, applied in the written order, and setup key \code{base\_region\_bending\_angle\_deg} for the detection angle & none \\
Add an inlet or an outlet & the geometry's \code{[ports]} identifies the surfaces; setup table \code{[[ports]]} selects each one's role, and the matrix row supplies its velocity and profile variables. Creating ports on a saved simulation stays refused while its existing port indices are unknown & none \\
\bottomrule
\end{tabularx}}
\guisource{Boundary conditions}
\end{frame}

\begin{frame}{Boundary conditions: what a key takes (3/3)}
\relax
{\ftstablesize\renewcommand{\arraystretch}{1.05}
\begin{tabularx}{\linewidth}{@{}JKZ@{}}
\guihead
Delete an existing inlet or outlet & setup keys \code{delete\_inlets} and \code{delete\_outlets}, port indices in the written order & none \\
Set the edge bluntness angles and the vertex merge tolerance & setup keys \code{geometric\_edge\_bluntness\_angle\_deg} and \code{vertex\_merge\_tolerance\_m}; the trailing-edge bluntness angle through setup table \code{[[flags]]}, which gives a one-value command a word of your own, or setup table \code{[[raw]]} & none \\
\bottomrule
\end{tabularx}}
\par\vspace{1mm}
\begin{alertblock}{The trailing edge by file is the default route}
\small Detection is an angle criterion, and a twisted blade's trailing edge is not a crease. The file names exactly the edges; the run checks the solver's count of imported edges against it.
\end{alertblock}
\guisource{Boundary conditions}
\end{frame}
